\newpage
\section[Theories, Methods and Current Research on Emotions in Library and Information Science, Information Retrieval and Human-Computer Interaction (53)]{Theories, Methods and Current Research on Emotions in Library and Information Science, Information Retrieval and Human-Computer Interaction (53) \cite{Lopatovska2011Emotions}}

\begin{enumerate}
    \item \textit{Category:} This paper is a literature review that examines theories of emotions, methods for studying emotions, and their applications in Library and Information Science (LIS), Information Retrieval (IR), and Human-Computer Interaction (HCI). It does not present new experimental findings but provides an overview and synthesis of existing research. The focus is on discussing how emotions are conceptualized and measured in these fields and how they influence human-computer interactions. 
    \item \textit{Context:} The paper is related to research in psychology, neuroscience, and affective computing, drawing from cognitive and somatic theories of emotions. It discusses the discrete and dimensional models of emotion representation and their relevance to human-computer interactions. The authors reference influential researchers whose work is foundational in understanding emotional expressions, physiological responses, and affective computing. The paper integrates findings from various disciplines to provide a comprehensive view of emotion research in information science.
    \item \textit{Correctness: }The assumptions made in the paper are well-grounded in prior research, as the authors review multiple emotion theories and evaluation methods without asserting a single dominant perspective. They acknowledge that emotion research is still evolving and that different models provide varying insights. The paper discusses methodological challenges, such as limitations in physiological measurements, self-reports, and observational techniques, and highlights the need for interdisciplinary collaboration to improve research accuracy. The balanced approach suggests that the paper’s claims are valid and credible.
    \item \textit{Contributions: }Contributions: The paper provides a comprehensive survey of emotion theories and methodologies used in Library and Information Science (LIS), Information Retrieval (IR), and Human-Computer Interaction (HCI). It organizes and compares different methods of studying emotions, including neurophysiological, observational, and self-report techniques, discussing their advantages and limitations. Additionally, it identifies gaps in interdisciplinary research, emphasizing the need for further collaboration between psychology, information science, and computing. The authors also suggest future research directions, particularly in the development of emotion-aware computing systems and their role in enhancing user experience.
    \item \textit{Clarity: }The paper is well-structured and systematically organized, making it accessible to researchers across multiple disciplines. It presents a thorough review of existing research while maintaining an objective and neutral stance. The paper is a valuable and informative resource, providing a broad overview of emotion research in computing and information science.
\end{enumerate}

    \textit{Refer to} \hyperref[fig3]{Appendix A3} for a figure on the \textsc{sources of facial expressions}
